\documentclass[aspectratio=169]{beamer}
\usetheme{default}
\setbeamertemplate{navigation symbols}{}
\usepackage{amsmath,amssymb}
\title{Kodaira--Spencer and the tangent space of $\mathrm{Def}(X)$}
\subtitle{Huybrechts, \S1}
\date{Reading seminar}
\begin{document}
\begin{frame}\titlepage\end{frame}
\begin{frame}{1. First-order variation}
Let $\pi:\mathcal X\to B$ be a smooth proper family with $X=\pi^{-1}(0)$.
The relative tangent sequence restricted to $X$ is
\[
0\to T_X\to T_{\mathcal X}|_X\to
T_0B\otimes\mathcal O_X\to0.
\]
Its connecting homomorphism is the Kodaira--Spencer map
\[
\rho:T_0B\to H^1(X,T_X).
\]
For a semiuniversal family $\rho$ is an isomorphism.
\end{frame}
\begin{frame}{2. Why a class in $H^1(T_X)$?}
Choose local trivialisations with changes of coordinates
$(x,t)\mapsto(\psi_{ij}(x,t),t)$.
For $v\in T_0B$, differentiating at $t=0$ gives
\[
\theta_{ij}=\left.\frac d{d\varepsilon}\right|_0
\psi_{ij}(x,\varepsilon v).
\]
These vector fields give a Cech $1$-cocycle;
changing trivialisations changes it by a coboundary.
Thus $\rho(v)=[\theta_{ij}]\in H^1(X,T_X)$.
\smallskip

{\scriptsize From HP \texttt{deformations.tex}, Kodaira--Spencer map;
\texttt{complex-geometry.tex}, relative tangent sequence.}
\end{frame}
\begin{frame}{3. The IHS computation}
Dolbeault gives $H^1(X,T_X)\simeq
H_{\bar\partial}^{0,1}(X,T^{1,0}X)$.
For an IHS manifold, contraction with the nondegenerate
holomorphic symplectic form $\sigma$ is an isomorphism
\[
T_X\xrightarrow{\ \iota_{(-)}\sigma\ }\Omega_X^1.
\]
Hence, by Dolbeault and the Kahler Hodge decomposition,
\[
\boxed{T_0\operatorname{Def}(X)
\simeq H^1(X,T_X)\simeq H^1(X,\Omega_X^1)
\simeq H^{1,1}(X).}
\]
The Bogomolov--Tian--Todorov theorem ensures that
$\operatorname{Def}(X)$ is smooth.
\end{frame}
\end{document}
